\color{unidarkgreen}{\subsection{Защита от замыкания на землю в сети НН реактора (контроль изоляции)}}
\color{black}

\begin{enumerate}[label=\arabic{section}.\arabic{subsection}.\arabic*,  align=left, labelsep=4pt, leftmargin=0pt, itemindent=57pt]

\item
Устройство контроля изоляции, представляет собой орган (функцию) максимального напряжения
нулевой последовательности устройства защиты КО, подключенного к цепям ТН на стороне НН УШР и
действующий на сигнал с заданной выдержкой времени.

\item
В связи с тем, что ошиновка НН УШР имеет достаточно локальный характер (как правило,
ограничивается собственно шинным мостом с присоединением не более двух ячеек КРУ ТМП), по
условиям эксплуатации и согласно требованиям безопасности персонала, защита от замыкания на
землю в сети НН реактора может быть введена с действием на отключение УШР используя программную
накладку \textbf{«Откл. КИ НН»}.

\item
Значение срабатывания по напряжению нулевой последовательности задается параметром
\textbf{«3U0ср КИ НН»}. Для исключения ложного пуска защиты при возникновении неисправности цепей
напряжения предусмотрен блокирующий орган, включенный на напряжение обратной
последовательности. Значение блокировки по напряжению обратной последовательности задается
параметром «\textbf{U2бл КИ НН»}. Действие на срабатывание формируется с задержкой по времени
\textbf{«Тср КИ НН»}.

\begin{figure}[H]
\centering
\includegraphics[width=0.7\textwidth,height=0.7\textheight,keepaspectratio]{\fbpath/КИ НН/img/KI_NN.pdf}
\captionof{figure}{Функциональный блок КИ НН}\label{pict_kinn:KI_NN}
\end{figure}

\small
\begin{longtable}{|>{\centering\arraybackslash}m{5.3cm}|>{\centering\arraybackslash}m{4.2cm}|>{\centering\arraybackslash}m{1cm}|>{\centering\arraybackslash}m{3.3cm}|>{\centering\arraybackslash}m{1.8cm}|}
% Первая страница: полный заголовок
\caption{Перечень уставочных параметров КИ НН\hfill\vspace{-0.5\baselineskip}}\label{kinn:tbl1}\\
\hline
\rowcolor{gray!30}
Наименование уставки & Диапазон & Шаг & Значение по умолчанию & Единица измерения \\
\hline
\endfirsthead

% Промежуточные страницы: «Продолжение таблицы»
\caption*{\hspace{3pt}\emph{Продолжение таблицы \ref{kinn:tbl1}}\hfill\ }\vspace{-0.5\baselineskip} \\
\hline
\rowcolor{gray!30}
Наименование уставки & Диапазон & Шаг & Значение по умолчанию & Единица измерения \\
\hline
\endhead

%===m>КИ НН|-
%===m
\end{longtable}
\normalsize

\end{enumerate}